\chapter*{Appendix H: Homotopy-Theoretic Tools for Semantic Structures}
\addcontentsline{toc}{chapter}{Appendix H: Homotopy-Theoretic Tools for Semantic Structures}

\section{Paths and Homotopies}

\subsection{Semantic Path Spaces}

Let $\mathcal{M}$ be a semantic manifold.  
A path is a continuous map
\[
\gamma : [0,1] \to \mathcal{M}.
\]

\begin{definition}[Path Space]
The path space of $\mathcal{M}$ is
\[
P(\mathcal{M}) = \{ \gamma:[0,1]\to\mathcal{M} \mid \gamma \text{ continuous} \}.
\]
\end{definition}

\subsection{Homotopies}

A homotopy between $f,g : X\to\mathcal{M}$ is a continuous
\[
H: X\times[0,1]\to\mathcal{M}
\]
satisfying $H(x,0)=f(x)$ and $H(x,1)=g(x)$.

\section{Homotopy Equivalence}

\begin{definition}[Homotopy Equivalence]
Spaces $X$ and $Y$ are homotopy equivalent if there exist maps
\[
f:X\to Y, \quad g:Y\to X
\]
such that $g\circ f\simeq \mathrm{id}_X$ and $f\circ g\simeq \mathrm{id}_Y$.
\end{definition}

\begin{definition}[Semantic Homotopy Equivalence]
Two semantic manifolds are homotopy equivalent if the equivalence preserves RSVP stratification and curvature strata.
\end{definition}

\section{Fibrations and Stratified Fibrations}

\subsection{Fibrations}

A map $p:E\to B$ is a fibration if for any space $X$ and any homotopy $H:X\times[0,1]\to B$, and any lift of $H(\cdot,0)$ to $E$, there exists a lift of $H$ to $E$.

\subsection{Stratified Fibrations}

A map $\pi:\mathcal{M}\to \mathcal{S}$ between stratified semantic manifolds is a stratified fibration if it is a fibration on each stratum and preserves codimension.

\section{Homotopy Limits and Colimits}

\subsection{Homotopy Colimit}

Given a diagram $D:\mathcal{I}\to \mathbf{Top}$,
\[
\operatorname*{hocolim}_{i\in\mathcal{I}} D(i)
\]
is constructed via the thick geometric realization of the simplicial replacement of $D$.

\subsection{Homotopy Limit}

Dually, the homotopy limit
\[
\operatorname*{holim}_{i\in\mathcal{I}} D(i)
\]
is the mapping space into the cosimplicial replacement.

\section{Applications to Semantic Structures}

\subsection{Hypergraph Realizations}

For a Polyxan hypergraph $\mathcal{G}$, its geometric realization $|\mathcal{G}|$ is assembled from simplices associated to cells.

\begin{definition}[Homotopy Type of a Hypergraph]
The homotopy type of $\mathcal{G}$ is the homotopy type of $|\mathcal{G}|$.
\end{definition}

\subsection{Embeddings}

Given an embedding $\iota:\mathrm{Inc}(\mathcal{G})\to\mathcal{M}$, the induced map  
\[
|\iota|:|\mathcal{G}|\to\mathcal{M}
\]
preserves homotopy classes.

\section{Postnikov Towers and Semantic Decomposition}

\subsection{Postnikov Sections}

Let $X$ be a semantic manifold.  
Its $n$-th Postnikov section $P_n X$ satisfies:
\[
\pi_k(P_n X) = \pi_k(X) \quad (k\le n),
\qquad
\pi_k(P_n X) = 0 \quad (k>n).
\]

\subsection{Semantic Decomposition}

A semantic manifold decomposes via a tower
\[
X \to \cdots \to P_nX \to P_{n-1}X \to \cdots \to P_1X.
\]

\section{Homotopy Groups}

\subsection{Definition}

For $x\in X$,
\[
\pi_n(X,x)
\]
is the set of homotopy classes of maps $(S^n,*)\to(X,x)$.

\subsection{Semantic Homotopy Groups}

\begin{definition}[Semantic Homotopy Group]
$\pi_n^{\mathrm{sem}}(\mathcal{M})$ is the subset of $\pi_n(\mathcal{M})$ represented by loops or spheres that remain within a fixed curvature stratum.
\end{definition}

\section{Homotopy Pushouts and Model Merging}

\subsection{Pushouts}

Given $X\leftarrow A\to Y$, the homotopy pushout is
\[
X \cup^{h}_A Y.
\]

\subsection{Semantic Pushouts}

For embeddings of hypergraph incidence complexes,
\[
|\mathcal{G}_1| \cup^{h}_{|\mathcal{H}|} |\mathcal{G}_2|
\]
models the geometric merging of semantic structures.

\section{Mapping Spaces and Function Complexes}

\subsection{Mapping Space}

The mapping space is
\[
\mathrm{Map}(X,Y)
\]
with compact-open topology.

\subsection{Semantic Mapping Space}

\[
\mathrm{Map}_{\mathrm{sem}}(X,Y)
  = \{ f:X\to Y \mid f \text{ preserves curvature strata} \}.
\]

\section{Classifying Spaces}

\subsection{Nerve Construction}

Given a category $\mathcal{C}$, its nerve $N\mathcal{C}$ is a simplicial set.

\subsection{Classifying Space}

The geometric realization $B\mathcal{C}=|N\mathcal{C}|$ is the classifying space of $\mathcal{C}$.

\section{Homotopy Coherence}

\subsection{Coherent Diagrams}

A diagram $D:\mathcal{I}\to\mathbf{Top}$ is homotopy coherent if all compositions commute up to specified homotopies satisfying higher coherence laws.

\subsection{Semantic Coherence}

Semantic diagrams (flows, embedding towers, reconciliation sequences) are homotopy coherent when each coherence cell preserves curvature and entropy.

\section{Homotopy of Semantic Flows}

\subsection{Flow Homotopy}

Two RSVP flows $\Phi_t$ and $\Psi_t$ are homotopic if there exists
\[
H: \mathcal{M}\times[0,1]_t \times[0,1]_s \to \mathcal{M}
\]
preserving stratification.

\subsection{Homotopy Invariance}

Energy descent, quine-stability, and curvature monotonicity are homotopy invariants for semantic flows.

\section{Homotopy Characterisation of the Universal Quine}

\subsection{Terminal Homotopy Type}

The universal media quine $\mathfrak{U}$ satisfies
\[
\pi_n^{\mathrm{sem}}(\mathfrak{U}) = 0
\quad \forall n\ge 1.
\]

\subsection{Uniqueness}

Any semantic manifold whose homotopy groups vanish within each curvature stratum is homotopy equivalent to $\mathfrak{U}$.


